% !TeX root = ../memoire.tex

\chapter*{Conclusion}
\addcontentsline{toc}{chapter}{Conclusion}
\markboth{Conclusion}{Conclusion}

La conservation des Archives historiques dans leur lieu d'origine et leur ouverture à des espaces documentaires plus larges ne constituent pas deux orientations contradictoires. L'étude menée à partir du cas de la Faculté de médecine de Montpellier montre qu'elles relèvent d'une même exigence : préserver l'intégrité et l'intelligibilité des fonds tout en rendant possibles leur repérage, leur consultation et leur réutilisation. La question posée par le projet de diffusion numérique ne pouvait être résolue par le seul choix d'une interface. Elle engage l'histoire et l'organisation des archives, leur statut juridique, la production de leurs descriptions, leur \gls{signalement} et les relations permettant à leurs différentes représentations de circuler au-delà de leur lieu de conservation.

L'histoire des Archives historiques éclaire la situation actuelle des fonds. Leur circulation entre officiers, professeurs et autorités, les restitutions du XVI\textsuperscript{e} siècle, les inventaires successifs, les reclassements du XVIII\textsuperscript{e} siècle puis l'intervention de Joseph Calmette au début du XX\textsuperscript{e} siècle ont laissé leur empreinte sur leur composition et leur organisation. Les cotes et les instruments de recherche héritent de ces opérations. Leur reprise suppose de comprendre les états antérieurs du classement et les déplacements de documents afin que l'actualisation des descriptions conserve les informations nécessaires à l'intelligibilité des fonds.

Cette connaissance prend une portée juridique et institutionnelle particulière. Dans la perspective de la demande de dérogation au versement, l'Université doit garantir durablement la conservation des fonds, leur classement, leur communicabilité et leur accès. La campagne engagée depuis 2023 répond à cette exigence par une progression qui doit rester centrée sur les archives. La diffusion numérique peut accompagner ce travail et en élargir les effets, mais elle dépend du traitement archivistique préalable.

Les expérimentations conduites pendant le stage ont montré que cette diffusion ne pouvait être ramenée à la création du mini-site initialement envisagé. La production des descriptions, leur \gls{signalement}, la conservation des reproductions et leur consultation peuvent relever d'environnements distincts, à condition que les relations entre le document, sa cote, sa description et sa reproduction soient maintenues. L'\gls{architecture-documentaire} distribuée proposée à l'issue de l'étude dissocie les fonctions tout en maintenant les liens nécessaires à l'intelligibilité des ressources. Elle autorise une progression dans laquelle le traitement et le \gls{signalement} des fonds peuvent avancer sans attendre la stabilisation de l'ensemble des choix techniques.

Cette organisation ouvre les Archives historiques à des ensembles documentaires qui dépassent leur périmètre. Les collections bibliographiques, muséales et scientifiques de l'Université, comme les sources conservées dans d'autres établissements, documentent parfois les mêmes personnes, institutions, enseignements ou pratiques sans relever des mêmes modèles descriptifs. Les \glspl{referentiel}, les \glspl{identifiant-perenne} et les réseaux de \gls{signalement} permettent d'établir ces relations sans réunir les ressources dans un système unique. La candidature \emph{Mémoire du monde} en fournit une illustration à l'échelle patrimoniale en définissant un ensemble cohérent à partir de documents répartis entre plusieurs institutions. Certaines sources peuvent être sélectionnées et transformées pour constituer des \glspl{corpus-numerique}, des transcriptions ou des données structurées répondant à des questions de recherche particulières. À mesure que s'ajoutent ces représentations, la \gls{provenance-donnees} devient essentielle pour maintenir leur relation avec les documents et les contextes archivistiques dont elles sont issues.

Le travail demeure celui d'une étude préalable. Le classement des archives modernes et contemporaines se poursuit, certains instruments restent à stabiliser et plusieurs décisions dépendent encore d'arbitrages institutionnels ou techniques. Les chaînes de diffusion devront être éprouvées en production, tandis que les rapprochements entre collections et la constitution de \glspl{corpus-numerique} devront partir d'ensembles précisément identifiés. Ces incertitudes n'empêchent pas de poursuivre les opérations qui peuvent être stabilisées indépendamment des outils : connaissance et traitement des fonds, harmonisation des descriptions, \gls{signalement}, attribution d'\glspl{identifiant-perenne} et documentation des relations entre les ressources. Cette progression limite la dépendance à un environnement technique particulier et prépare les données à des usages qui pourront évoluer avec les choix de l'établissement.

La demande de dérogation prévue à l'horizon 2029 donne à cette progression un cadre concret. Le maintien des archives dans le bâtiment historique reposera sur la capacité de l'Université à démontrer qu'elle en assure durablement la conservation, l'intelligibilité et l'accessibilité ; leur ouverture peut, dans le même temps, les inscrire dans des espaces documentaires et scientifiques qui dépassent largement la Faculté. Les fonds demeurent dans leur lieu de conservation tandis que leurs descriptions circulent dans des réseaux, que leurs reproductions peuvent être consultées ailleurs, que les entités qu'ils documentent sont reliées à d'autres collections et que certaines séries deviennent les sources de nouveaux \glspl{corpus-numerique}. Conserver à demeure ne signifie donc pas maintenir les archives à l'écart, mais organiser ces circulations sans perdre le lien avec les documents, leur contexte et l'institution qui en assume la responsabilité. C'est dans cette articulation que prend pleinement sens le titre de ce mémoire : \emph{À demeure, en partage}.